
\subsubsection{Data Attribution Methods}

Training data attribution has evolved from influence functions \citep{koh2017understanding} through increasingly scalable approximations. TracIn \citep{pruthi2020tracin} introduced gradient tracing across checkpoints. TRAK \citep{park2023trak} achieved speedup through random projection of gradients. Kronfluence \citep{grosse2023kronfluence} applied Kronecker-factored curvature approximation (K-FAC) to enable influence computation at scale.

These methods represent different computational strategies, but prior work has not systematically characterized what each strategy measures. This work addresses this gap by showing that computational differences manifest as systematic differences in sensitivity to influence modes.

\subsubsection{Attribution Method Fragility and Benchmarking}

Basu et al. (2020) demonstrated that influence functions are fragile in deep networks, with approximation errors increasing with network depth. This finding raised concerns about gradient-based attribution reliability. The present work suggests that apparent fragility may be reframed: rather than random errors, different methods may exhibit systematic sensitivity patterns that constitute predictable fingerprints.

The DATE-LM benchmark \citep{jiao2025datelm} provided the first unified evaluation of LLM attribution methods, revealing that no single method dominates across tasks. This work complements DATE-LM by investigating why methods differ: the inductive biases embedded in each algorithm's mathematics may create different sensitivities to memorization, feature transfer, and spurious association.

Several efficiency-focused works have pushed attribution to larger scales. LoRIF \citep{li2026lorif} achieved $20\times$ storage reduction on 70B models, and GraSS \citep{hu2025grass} demonstrated 165\% throughput improvement through gradient sparsification. GGDA \citep{ley2024ggda} introduced group-level attribution with 10-$50\times$ speedup. These works focus on computational efficiency; this work focuses on characterizing what methods measure.

\subsubsection{Learning Dynamics and Influence Modes}

Research on learning dynamics has identified distinct phases and patterns in neural network training. Studies of memorization \citep{feldman2020memorization} show that models memorize atypical examples to achieve low training loss. Work on spurious correlations \citep{sagawa2020spurious} reveals that models exploit dataset shortcuts.

These distinct influence modes---memorization, feature transfer, spurious association---have been studied in isolation. The contrastive mode probing framework provides a unified lens for measuring how attribution methods respond to each mode.
